% !TEX program = xelatex
\documentclass[UTF8,a4paper,11pt]{ctexart}
\usepackage[margin=2.2cm]{geometry}
\usepackage{graphicx}
\usepackage{booktabs}
\usepackage{enumitem}
\usepackage{xcolor}
\usepackage{caption}
\usepackage{tabularx}
\usepackage[colorlinks=true,linkcolor=blue!55!black,urlcolor=blue!55!black]{hyperref}
\setlist{itemsep=2pt, topsep=4pt}
\captionsetup{font=small, labelsep=quad, skip=4pt}
\graphicspath{{img-zh/}}
\newfontfamily\keysfont{Apple Symbols}
\newcommand{\cmdkey}{{\keysfont ⌘}}
\newcommand{\optkey}{{\keysfont ⌥}}
\newcommand{\shiftkey}{{\keysfont ⇧}}
\newcommand{\swapsym}{{\keysfont ⇄}}
\newcommand{\ui}[1]{「#1」}
% 插图：浮动排版，可选宽度（行宽的倍数，默认 0.7），标签 fig:文件名。
\newcommand{\shot}[3][0.7]{\begin{figure}[!htbp]\centering\includegraphics[width=#1\linewidth]{#2}\caption{#3}\label{fig:#2}\end{figure}}
\newcommand{\fig}[1]{图~\ref{fig:#1}}
\renewcommand{\topfraction}{0.9}
\renewcommand{\bottomfraction}{0.8}
\renewcommand{\textfraction}{0.07}
\renewcommand{\floatpagefraction}{0.8}
\setcounter{topnumber}{3}
\setcounter{totalnumber}{4}
\ctexset{section={format=\Large\bfseries}}
% 带圈数字用中文字体排版。
\xeCJKsetcharclass{"2460}{"2473}{1}

\title{\textbf{BiWrite 使用说明}}
\author{版本 0.2.2}
\date{2026 年 10 月}

\setlength{\emergencystretch}{2em}
\begin{document}
\maketitle
\thispagestyle{empty}

\begin{center}
开源仓库 \url{https://github.com/zzccppp/biwrite}\\[2pt]
写作技能 research-builder \url{https://github.com/qzkinhit/research-builder}
\end{center}

\begin{abstract}
BiWrite 是一个中英双语的论文写作编辑器。左侧编写正文，右侧逐段显示另一种语言，只有改动过的段落会重新翻译。LaTeX 论文可以在软件内编译成 PDF，点击 PDF 中的句子即可在源码中选中该句。写作助手按 research-builder 写作技能的规则润色、修改、回答问题、参照已有论文仿写和生成图表。第~\ref{sec:flow}~节按一篇论文从新建到投稿的顺序介绍完整的写作流程，之后各节按功能分步介绍。截图中的红框标出要操作的按钮或要看的区域，正文中的 ①②③ 与红框旁的编号一一对应。
\end{abstract}

\section{功能一览}
\label{sec:overview}
点击下面的功能名称可以跳到对应的章节。
\begin{description}[leftmargin=0pt, style=nextline]
\item[{\hyperref[sec:flow]{完整写作流程}}] 从用模板新建论文开始，经过写作、仿写、插入图表、编译 PDF、改中文，到导出中文版，十个步骤串起全部功能。
\item[{\hyperref[sec:write]{双语写作}}] 左侧写英文，右侧逐段显示中文，只有改动过的段落会重新翻译。顶部的切换按钮可以改为直接编辑中文，英文随之修改，没改的段落保持原文。提前切换时，尚未译完的段落在后台译好后自动替换到左侧。
\item[{\hyperref[sec:zhfile]{中文原文的论文}}] 打开中文文件时自动改为中译英，保存时写入的始终是文件本身的语言。原文语言判断错误时一键改正。
\item[{\hyperref[sec:latex]{LaTeX 编译与 PDF 定位}}] 在软件内编译 PDF。点击 PDF 中的句子，源码中对应的句子被选中，并可直接润色。光标所在段落可以反向定位到 PDF，编译错误点击即跳到源码行。
\item[{\hyperref[sec:zhpdf]{中文 PDF}}] 英文译完后一键编译中文 PDF，并可导出中文 PDF 和可以单独编译的中文 .tex 文件。
\item[{\hyperref[sec:pair]{导入已有中文译文}}] 与 \texttt{paper\_zh.tex} 或 \texttt{sections\_zh/} 逐段配对。改任何一侧，另一侧文件中对应的段落原位更新，其余内容不变。
\item[{\hyperref[sec:assist]{写作助手}}] 润色、按指令修改、提问、参照已有论文仿写、生成图表、按译文改原文。参照论文可以直接添加 \texttt{.tex} 文件，生成的图表缺少宏包时插入时自动补上。任务在后台运行，结果由你决定是否采纳。
\item[{\hyperref[sec:templates]{论文模板}}] 内置 VLDB、ICLR 2027 和 IEEE TII 模板，可以导入自己的项目作为模板，也可以导出当前项目。
\item[{\hyperref[sec:provider]{多种接口与号池}}] 支持 OpenAI 兼容接口、OpenAI Responses 接口和 Anthropic 接口。一个接口可以挂多个密钥，均匀分配请求，限流时自动换号重试。
\item[{\hyperref[sec:log]{请求日志}与\hyperref[sec:update]{版本更新}}] 日志记录每次请求的模型、推理强度和服务等级。软件内可以安装 GitHub 上发布的任意版本。
\item[{\hyperref[sec:tray]{关闭到托盘}}] 关闭窗口后 BiWrite 留在菜单栏或任务栏通知区域，文档和翻译进度都保留，点图标即可回到窗口。
\end{description}
\fig{o-editor} 是主界面，①为源文编辑区，②为译文区。

\shot[0.8]{o-editor}{主界面：①左侧编辑英文 LaTeX，②右侧逐段显示中文}

\section{完整写作流程}
\label{sec:flow}
本节按一篇论文从新建到投稿的顺序，把各项功能串成十个步骤。每一步给出操作要点，细节见后文对应的章节。

\subsection{第 1 步：用模板新建论文}
点工具栏的\ui{新建}，在\ui{论文与模板}窗口中找到要投的会议或期刊（\fig{f-new} ①），点\ui{用此模板新建…}（②），再为论文起一个文件夹名。软件复制模板文件并打开主文件。导入自己的模板见第~\ref{sec:templates}~节。

\shot{f-new}{①ICLR 2027 模板，②用此模板新建}

\subsection{第 2 步：写作与实时翻译}
在左侧写英文，停止输入 0.8 秒后，右侧显示改动过的段落的中文（\fig{o-editor}）。也可以先写中文初稿，软件识别出原文是中文后，会把中文翻译成英文，见第~\ref{sec:zhfile}~节。

\subsection{第 3 步：参照已有论文仿写}
写引言、相关工作这类结构相对固定的段落时，可以让写作助手参照一篇已发表的论文仿写。
\begin{enumerate}
\item 把光标放在新段落要接在后面的那一段，按 \cmdkey J 打开写作助手，选\ui{仿写}（\fig{f-imitate} ①）。②显示新段落将插在哪一段之后。光标放在空行上时，新段落插在光标处。
\item 在③写明要求，包括写什么、写几段、写多长，例如\ui{仿照参照论文的引言，写两段说明图上下文学习的动机}。
\item 点\ui{参照样例}一行的\ui{+ 文件}（④），选择要模仿的论文，例如 UniClean 论文的 \texttt{.tex} 文件。LaTeX 文件只取正文并去掉注释，篇幅过长时截取前 12 万字。添加后标签上显示文件名和字数（⑤）。也可以用\ui{+ 文本}粘贴一段文字。
\item 点\ui{开始}（⑥）。
\end{enumerate}
模型模仿参照论文的段落安排、句式和论证方式，内容来自你的要求和论文的上下文，不照搬参照论文的句子、结论、数字和引用。论文中没有的实验结果，模型会写成 \texttt{[result]} 这样的占位标记，留给你补上。

\shot{f-imitate}{①仿写，②新段落的位置，③写作要求，④添加参照论文，⑤已添加的论文，⑥开始}

结果卡片（\fig{f-written}）显示新写的段落（①）、逐段对照的中文（②）和写法说明（③）。点\ui{插入}（④），新段落插在目标段落之后，右侧直接显示卡片中的中文，不再重新请求翻译（\fig{f-inserted}）。

\shot[0.55]{f-written}{①新写的段落，②对照中文，③写法说明，④插入}
\shot[0.8]{f-inserted}{①插入左侧的新段落，②右侧直接显示对应的中文}

\subsection{第 4 步：润色、修改与提问}
写好的段落用\ui{润色}统一措辞，用\ui{按指令改}按要求调整，用\ui{提问}检查论证是否成立。操作方法见第~\ref{sec:assist}~节。

\subsection{第 5 步：插入图表}
在写作助手中选\ui{图表}，在要求中写清图表的内容，例如\ui{四个基准上的准确率，做成表格}。结果卡片给出可以直接插入的 LaTeX 代码（\fig{f-table} ①）。生成代码的请求中列出文档已经加载的宏包，并要求优先使用这些宏包。代码还需要文档没有加载的宏包时，卡片会列出这些宏包（②）。当前文件就是主文件时，点\ui{插入}（③）会把代码插在段落之后，同时把对应的 \texttt{\textbackslash usepackage} 加到 \texttt{\textbackslash begin\{document\}} 之前（\fig{f-table-in}）。当前文件是某一节的分文件时，卡片提示要在主文件导言区加入哪几行。

\shot{f-table}{①表格代码，②需要补充的宏包，③插入}
\shot{f-table-in}{①插入时自动加入导言区的宏包}

仿照已有的图画图时，选\ui{图表}（\fig{f-figref} ①），在要求中写明要画什么、数据是什么（②），再用\ui{参照样例}一行的\ui{+ 图片}（③）添加那张图，添加后显示为带缩略图的标签（④），最后点\ui{开始}（⑤）。模型按参考图的图表类型、版式、配色、标记、线型、字体和图例位置，用 TikZ 或 pgfplots 画出同样风格的图，数据和标签来自你的要求和论文，不照抄参考图上的数字。参考图片支持 PNG、JPEG、WebP 和 GIF，每张不超过 8 MB，所用的模型需要支持图片输入。

\shot{f-figref}{仿照参考图画图：①图表，②要求，③添加参考图片，④已添加的图片，⑤开始}

\subsection{第 6 步：编译 PDF 并在 PDF 上修改}
切到右侧的 PDF 页签，按 \cmdkey B 编译。在 PDF 上点击一个句子，源码中对应的句子被选中，可以直接润色或修改，见第~\ref{sec:latex}~节。投稿版论文页边的行号不影响定位。

\subsection{第 7 步：切换为改中文}
点\ui{EN \swapsym{} 中}，左侧改为中文，按中文修改，英文随之更新，见第~\ref{sec:write}~节。只想改某一段时，点右侧这一段的铅笔按钮，用\ui{按译文改}改写这段中文，英文按你的中文修改，见第~\ref{sec:pair}~节。

\subsection{第 8 步：编译和导出中文版}
在 PDF 页签切到\ui{中文}，编译中文 PDF，再用\ui{另存 PDF…}和\ui{导出中文 .tex…}保存中文 PDF 和中文 \texttt{.tex}，见第~\ref{sec:zhpdf}~节。

\subsection{第 9 步：与中文版同步维护}
把导出的中文 \texttt{.tex} 放在英文主文件旁边，命名为 \texttt{论文名\_zh.tex}。下次打开英文文件时两者自动配对，此后改任意一侧，另一侧文件中对应的段落原位更新，见第~\ref{sec:pair}~节。

\subsection{第 10 步：保存与归档}
\cmdkey S 保存，\ui{另存为}保留不同阶段的版本。论文定稿后，可以在\ui{新建}窗口用\ui{导出当前项目…}把它打包成模板，供下一篇论文使用。

\section{日常写作}
\label{sec:write}

\subsection{打开与编辑}
点工具栏的\ui{打开}（\fig{w-open} ①，快捷键 \cmdkey O）选择 \texttt{.tex}、\texttt{.md} 或 \texttt{.txt} 文件。多文件的 LaTeX 项目会在标题栏显示文件列表（②），可以在项目文件之间切换。右侧每个方块对应源文的一个段落、标题或图注，点击方块，光标跳到对应段落，当前段落在两侧同时高亮（③）。

停止输入 0.8 秒后，只有改动过的段落会重新翻译，并以流式方式显示。公式、引用、交叉引用、标签、网址和注释在发给模型前被替换成占位符，返回后按原样还原。

\shot[0.8]{w-open}{①打开文件，②项目文件列表，③当前段落的译文}

\subsection{中英文切换修改}
工具栏中间的\ui{EN \swapsym{} 中}按钮（\fig{w-swap} ①）决定左侧编辑哪种语言。
\begin{enumerate}
\item 默认左侧编辑英文，右侧显示中文。
\item 点切换按钮后，左侧变为中文（\fig{w-swapped} ②），可以直接改中文，右侧显示英文（③）。你改了哪一段中文，就只重新生成这一段的英文，并以原来的英文为基础做最小修改，其余段落保持你原来的英文措辞。
\item 再点一次（①）切回英文编辑。没有改过的段落会逐字还原。
\end{enumerate}
切换需要每个段落都已有译文。点击时如果还有段落正在翻译，按钮会显示等待标记，等这些段落译完后自动切换。不想等可以再点一次，立即切换。尚未译完的段落先保留原文，后台继续翻译，译好后自动替换到左侧（\fig{f-filling}）。替换前你改过的段落保持你的改动。切回英文时，这些段落的英文逐字还原。自动翻译暂停时不会在后台翻译，可以用\ui{全部}菜单中的\ui{继续翻译剩余段落}补上。

\shot{w-swap}{①切换编辑语言的按钮}
\shot[0.8]{w-swapped}{切换后：①按钮变为「中 \swapsym{} EN」，②左侧编辑中文，③右侧显示对应的英文}
\shot[0.8]{f-filling}{提前切换：①已切换为编辑中文，②尚未译完的段落正在后台翻译，③状态栏的提示}

保存时写入的始终是文件本身的语言。英文文件在编辑中文时保存，软件用当前译文组成英文后写入原文件。导入了已有的中文文件后，两个文件都会写入，见第~\ref{sec:pair}~节。

\subsection{全部菜单}
工具栏的\ui{全部}按钮（\fig{f-all} ①）打开一个菜单。
\begin{itemize}
\item ②\ui{重新翻译所有段落}：每段都重新请求翻译，不使用缓存。切换为编辑另一种语言时不可用，因为右侧是你文件中的原文。
\item ③\ui{继续翻译剩余段落}：只翻译还没有译文或翻译失败的段落，自动翻译暂停时也会执行。提前切换后，左侧仍是原文语言的段落也会译好后替换到左侧。
\item ④\ui{把原文当作中文，翻译成英文}（英文论文中显示为这一项，中文论文中显示为相反的一项）：原文语言判断错误时使用，见第~\ref{sec:zhfile}~节。
\end{itemize}

\shot{f-all}{①全部菜单，②重新翻译所有段落，③继续翻译剩余段落，④改变原文语言}

\subsection{中文原文的论文}
\label{sec:zhfile}
打开中文文件时，软件根据正文段落判断原文语言（只看正文，公式、引用和 LaTeX 命令不计在内）。中文明显占多数时，左侧编辑中文，右侧显示英文译文，状态栏显示 zh \textrightarrow{} en（\fig{f-zhfile}）。此时：
\begin{itemize}
\item 保存时写入的是左侧的中文。点切换按钮可以改为编辑英文，保存时软件用当前译文组成中文写回文件。
\item PDF 页签的\ui{中文}编译磁盘上的中文文件，\ui{英文}由英文译文编译。导出按钮变为\ui{导出英文 .tex…}。
\end{itemize}
中英文混合、判断不准时，可以手动改正。左侧文字明显是另一种语言时，点切换按钮会直接改为按这种语言翻译，不需要等待（\fig{f-retarget}）。其他情况在\ui{全部}菜单中选\ui{把原文当作中文，翻译成英文}或相反的一项。文件内容不变，只是按新的原文语言重新翻译。

\shot[0.8]{f-zhfile}{中文文件：①按钮为「中 \swapsym{} EN」且未切换，②右侧为英文译文，③状态栏显示 zh \textrightarrow{} en}
\shot[0.8]{f-retarget}{原文语言判断错误时点切换按钮：①直接改为编辑中文原文，②状态栏变为 zh \textrightarrow{} en，③提示}

\subsection{保存与另存为}
\ui{保存}（\fig{w-save} ①，\cmdkey S）覆盖当前文件。\ui{另存为}（②，\cmdkey\shiftkey S）以新文件名保存，默认建议一个不重名的名字，例如 \texttt{paper-2.tex}，原文件保持不变。

\shot{w-save}{①保存，②另存为}

\subsection{术语表}
工具栏的\ui{术语表}按钮打开术语表（\fig{w-glossary}）。在表格（①）中为术语指定中文译法，或勾选\ui{保留英文}。上方的按钮（②）可以导入和导出 CSV（两列 \texttt{term,translation}，UTF-8 编码）。修改术语后，只有提到这个术语的段落会重新翻译。

\shot[0.6]{w-glossary}{术语表：①术语与译法，②导入导出和保存}

\section{LaTeX 论文}
\label{sec:latex}

\subsection{打开 LaTeX 项目}
LaTeX 文件可以直接用\ui{打开}。也可以点\ui{新建}（\fig{l-files} ①），在弹出的窗口里选\ui{打开 LaTeX 文件夹…}，选择论文所在文件夹，软件会打开其中的主文件。多文件项目的标题栏显示文件列表（②），标有\ui{主文件}的是编译入口。LaTeX 文件的右侧窗格上方有\ui{译文}和\ui{PDF}两个页签（③）。

软件按以下顺序确定主文件：文件开头的 \texttt{\% !TEX root = …} 注释，然后是通过 \texttt{\textbackslash input}、\texttt{\textbackslash include} 链引入当前文件的主文件，嵌套多层也能找到。编译引擎依次参考 \texttt{\% !TEX program} 注释、项目中 \texttt{latexmkrc} 的 \texttt{\$pdf\_mode} 设置，以及导言区加载的宏包。

\shot[0.95]{l-files}{①新建与打开 LaTeX 文件夹，②项目文件列表，③译文与 PDF 页签}

\subsection{编译英文 PDF}
\begin{enumerate}
\item 点右侧的\ui{PDF}页签（\fig{l-compile} ①）。
\item 点\ui{编译}（②）或按 \cmdkey B（Windows 为 Ctrl+B）。未保存的改动会先保存，因为 PDF 来自磁盘上的文件。
\item 编译完成后显示 PDF。③处显示用时、错误数和警告数，④处的问题列表列出每条错误和警告，点击即跳到对应源码行。
\end{enumerate}
\ui{设置}中\ui{保存 .tex 文件后自动编译}默认开启，每次保存都会重新编译。文档有错误时仍会生成 PDF，问题列表中会列出这些错误。

\shot{l-compile}{①PDF 页签，②编译，③编译结果，④问题列表}

\subsection{点击 PDF 定位源码}
在 PDF 中单击一个句子（\fig{l-click} ①），左侧源码中对应的句子被选中（②）。点击位置旁会出现一个小菜单（③），提供\ui{润色}、\ui{按指令改…}和\ui{提问…}，可以直接对这句话发起写作助手任务。如果这句话在另一个项目文件里（例如 \texttt{sections/04\_method.tex}），软件会打开那个文件并选中句子。当前文件有未保存改动时，软件先提示，由你决定是否打开。

软件会核对点击处的文字。投稿版论文页边的行号、页眉这类叠加在正文上的内容会干扰 SyncTeX 的结果，这时软件按点击处的词语在当前文件和项目的其他文件中查找对应的句子。

\shot[0.8]{l-click}{①点击 PDF 中的句子，②源码中对应的句子被选中，③操作菜单}

\subsection{光标定位到 PDF}
把光标放在源码中的某一段（\fig{l-locate} ①），点 PDF 工具栏的定位按钮（②），或按 \cmdkey\optkey J（Windows 为 Ctrl+Alt+J），PDF 中对应的位置会高亮几秒（③）。

\shot[0.8]{l-locate}{①光标所在段落，②定位按钮，③PDF 中的对应位置}

\subsection{英文翻译完成后编译中文 PDF}
\label{sec:zhpdf}
中文 PDF 由当前译文组成，用 XeLaTeX 编译，并自动加入中文字体支持（ctex 宏包）。
\begin{enumerate}
\item 打开英文论文，等右侧译文全部完成。状态栏显示已翻译、进行中和等待中的段落数。
\item 在 PDF 页签点\ui{中文}（\fig{l-zh} ①）。第一次切换会自动编译，之后可以点\ui{编译}（②）重新生成。
\item 尚未翻译的段落会以英文排版，PDF 上方会提示这类段落的数量（③）。
\end{enumerate}
中文 PDF 和编译过程中的文件保存在论文文件夹下的隐藏文件夹 \texttt{.biwrite/zh/} 中，不会覆盖你的任何文件。

\shot{l-zh}{①切换到中文 PDF，②编译，③尚未翻译的段落提示}

\subsection{导出中文 PDF 与中文 .tex}
在中文 PDF 页面（\fig{l-export}）：
\begin{itemize}
\item ①\ui{另存 PDF…}保存一份中文 PDF，默认文件名为 \texttt{论文名.zh.pdf}。英文 PDF 用同一个按钮另存。
\item ②\ui{导出中文 .tex…}得到一个可以单独编译的中文 LaTeX 文件，默认名为 \texttt{论文名\_zh.tex}，用 XeLaTeX 编译即可。中文论文在这里导出英文版，按钮为\ui{导出英文 .tex…}，查看英文 PDF 时显示。
\item ③\ui{在文件夹中显示}在访达或资源管理器中显示当前 PDF。
\end{itemize}

\shot{l-export}{①另存 PDF，②导出中文 .tex，③在文件夹中显示}

\subsection{论文模板}
\label{sec:templates}
点\ui{新建}打开\ui{论文与模板}窗口（\fig{l-templates}）。
\begin{itemize}
\item 内置模板有 VLDB（PVLDB 第 20 卷，2027 年）、ICLR 2027 和 IEEE TII。前两个使用官方模板文件，TII 附带一个简短的论文骨架。
\item ①\ui{用此模板新建…}：为新论文起一个文件夹名，软件复制模板并打开主文件，切到 PDF 页签即可编译。
\item ②\ui{打开 LaTeX 文件夹…}：打开已有的论文项目。
\item ③\ui{导入文件夹…}和\ui{导入 .zip…}：把已有的论文项目保存为自己的模板，显示在\ui{我的模板}中。
\item ④\ui{导出当前项目…}：把当前论文项目打包成 zip，附带模板描述文件，可以发给他人导入。
\end{itemize}

\shot{l-templates}{①用模板新建论文，②打开 LaTeX 文件夹，③导入模板，④导出当前项目}

\section{导入已有的中文译文}
\label{sec:pair}
如果你已经有一份人工维护的中文版（例如 \texttt{paper.tex} 对应 \texttt{paper\_zh.tex}），BiWrite 可以直接使用它，不再调用接口翻译。两个文件配对后，改英文会更新中文文件中对应的段落，改中文会更新英文文件中对应的段落。

\subsection{自动配对}
打开英文文件时，软件按文件名查找对应的中文文件，找到并且段落能对上时自动配对。支持两类命名方式。
\begin{itemize}
\item 同一文件夹下加后缀：\texttt{paper.tex} 对应 \texttt{paper\_zh.tex}，后缀也可以是 \texttt{.zh}、\texttt{-zh}、\texttt{\_cn}。
\item 平行的语言文件夹：\texttt{sections\_en/04\_method.tex} 对应 \texttt{sections\_zh/04\_method.tex}，\texttt{en/} 对应 \texttt{zh/}。
\end{itemize}

\subsection{手动导入}
\begin{enumerate}
\item 打开并保存英文文件。
\item 点标题栏文件名右侧的\ui{导入译文}按钮（\fig{p-import} ①，链条图标）。
\item 选择中文文件。软件逐段对齐两个文件，标题栏出现中文文件名的标签（\fig{p-paired} ①），状态栏显示对上的段落数（②），右侧直接显示你的中文（③）。
\end{enumerate}
对齐按段落类型（正文、标题、图注）、共同出现的引用键、交叉引用、标签、公式、数字和专有名词，以及段落长度进行，图表在两个文件中位置不同也能对上。对不上的段落会单独翻译，但不会写入中文文件。

\shot{p-import}{①导入译文}
\shot[0.8]{p-paired}{配对后：①中文文件标签，②对上的段落数，③右侧显示你的中文}

\subsection{配对后的编辑与保存}
\begin{itemize}
\item 右侧直接显示你的中文，不产生任何请求。
\item 改一段英文，只有这一段的中文会以你原来的中文为基础重新修订。
\item 也可以直接改中文：点右侧段落上的铅笔按钮（\fig{p-edit} ①），写作助手打开\ui{按译文改}，输入框（②）里是这段的中文。改完点\ui{开始}（③），模型按你的中文最小幅度修改英文。结果卡片显示英文的改动（\fig{p-done} ①），点\ui{采纳}（②）写入英文，这段的中文就是你写的那句。
\item 保存时两个文件都会写入。中文文件里只有改动过的段落被原位替换，其余内容逐字节不变。如果某段译文还在生成，英文文件先保存，中文文件在译文完成后自动写入。
\item 点切换按钮后，左侧编辑的就是你的中文文件，改动以同样方式写回英文文件。
\item \ui{另存为}（\fig{p-saveas} ①）会把两个文件都按新名字另存，例如 \texttt{paper-2.tex} 和 \texttt{paper\_zh-2.tex}，原来的两个文件不变。
\end{itemize}

\shot[0.55]{p-edit}{①铅笔按钮，②改写这段的中文，③开始}
\shot[0.55]{p-done}{①英文的改动，②采纳}
\shot{p-saveas}{①另存为，②取消配对}

\subsection{配对后的 PDF 与取消配对}
配对后，PDF 页签中的\ui{中文}直接编译你自己的中文主文件（例如 \texttt{paper\_zh.tex}），而不是由译文临时组成的文档。点击中文 PDF 中的句子，会选中英文源码中对应的段落。点中文文件标签上的叉号（\fig{p-saveas} ②）即取消配对，之后保存不再写入中文文件。

\section{写作助手}
\label{sec:assist}

\subsection{打开与六种操作}
点工具栏\ui{助手}或按 \cmdkey J 打开右侧的写作助手（\fig{a-dock}）。它作用于选中的文字，没有选中时作用于光标所在段落（②）。
\begin{enumerate}
\item 在①选择操作。
\begin{description}[leftmargin=1em, style=nextline]
\item[润色] 按写作规则修改措辞，去除 AI 腔、拼接句子的破折号和分号，保持论断、数字、引用和公式不变。\cmdkey\shiftkey P 直接润色光标所在段落。
\item[按指令改] 按你写的要求修改，例如\ui{再简洁一些，写明基准测试的个数}。
\item[提问] 就选中的文字或整篇论文提问，回答下方可以继续追问。
\item[仿写] 按要求写新的段落，插在段落之后。添加参照论文后，模仿它的写法，见第~\ref{sec:flow}~节第 3 步。
\item[图表] 按描述生成 LaTeX 图或表，插入在段落之后，缺少的宏包一并补上，见第~\ref{sec:flow}~节第 5 步。
\item[按译文改] 输入框里预填这一段的译文，改写译文后运行，模型据此修改原文，见第~\ref{sec:pair}~节。
\end{description}
\item 在③写要求或问题（润色可以不写）。
\item ④\ui{上下文}决定发给模型的范围，可选仅目标、含前后段或全文。
\item ⑤\ui{参照样例}可以添加想模仿其写法的文字片段（\ui{+ 文本}）、整篇论文的 \texttt{.tex}、\texttt{.md} 或 \texttt{.txt} 文件（\ui{+ 文件}），或添加图片（\ui{+ 图片}，例如想参照其风格的图）。
\item 点\ui{开始}（⑥）。任务在后台运行，你可以继续写作。
\end{enumerate}

\shot{a-dock}{①操作，②作用对象，③要求，④上下文，⑤参照样例，⑥开始}

\subsection{确认结果}
任务完成后，结果卡片（\fig{a-result}）显示修改前后的差异（①，删除部分划线，新增部分高亮）、修订稿的对照译文（②）和中英文两份改动说明（③）。点\ui{采纳}（④）把修订写入当前位置，即使你在等待期间继续输入，软件也会找到原文的新位置。如果目标文字本身在等待期间被改过，卡片会提示，点\ui{让模型重新套用}即可。\ui{重新生成}再运行一次，\ui{放弃}丢弃结果，采纳后可以用 \cmdkey Z 撤销。

\shot{a-result}{①差异，②对照译文，③改动说明，④采纳}

\subsection{提问}
提问的回答显示在卡片中（\fig{a-ask} ①），下方的输入框（②）可以继续追问。仿写和图表的用法见第~\ref{sec:flow}~节第 3 步和第 5 步。

\shot{a-ask}{①回答，②追问}

\subsection{写作技能}
助手的规则来自开源写作技能 research-builder（MIT 协议）。软件自带一份，在\ui{设置}的\ui{写作技能}一节（\fig{a-skill}）可以①从 GitHub 更新到最新版本，②改用你自己的技能文件夹，③打开技能的开源仓库。

\shot{a-skill}{①从 GitHub 更新，②使用自己的技能文件夹，③技能仓库}

\section{请求日志}
\label{sec:log}
工具栏右侧的日志按钮（\cmdkey\shiftkey L）打开请求日志（\fig{g-log}）。①处可以按类型筛选、暂停记录、清空和打开日志文件。②的表格中每行是一次请求，列出发出的模型、推理强度和服务等级，以及服务端声明的实际值，不一致时会标出。点击一行展开详情（\fig{g-details} ①），包括 HTTP 状态、首字延迟、总耗时、输入与输出 token（含缓存和推理 token）、所用密钥的名称和末尾几位，以及重试和换密钥的说明。日志不记录正文和密钥。

\shot[0.8]{g-log}{①筛选与操作，②请求列表}
\shot[0.95]{g-details}{①一次请求的详情}

\section{版本更新}
\label{sec:update}
在\ui{设置}最下方的\ui{版本更新}点\ui{查看可用版本}（\fig{u-list} ①），列出 GitHub 上发布的全部版本及更新说明（②），包括旧版本和预发布的分支构建。选择版本后点\ui{安装}（③），软件下载对应的安装包，并按 GitHub 提供的 SHA-256 校验。macOS 上安装完成后点\ui{立即重启}（\fig{u-done} ①）进入新版本。Windows 上安装程序会在 BiWrite 退出后运行。软件启动时会检查是否有新版本，有则在状态栏提示，这一检查可以关闭。

\shot{u-list}{①查看可用版本，②版本列表，③安装}
\shot{u-done}{①安装完成后重启}

\section{安装}
\label{sec:install}

\subsection{macOS}
从 Releases 页面（\url{https://github.com/zzccppp/biwrite/releases}）下载 \texttt{BiWrite\_\textless 版本\textgreater\_universal.dmg}，它同时支持 Apple 芯片和 Intel 芯片的 Mac，要求 macOS 12 或更高版本。打开 dmg 后，把 BiWrite 拖进\ui{应用程序}文件夹。

安装包没有使用 Apple 签发的开发者证书签名，第一次打开时系统会拦截。放行步骤如下。
\begin{enumerate}
\item 双击 BiWrite，系统提示无法验证开发者，点\ui{完成}关闭提示。
\item 打开\ui{系统设置}，进入\ui{隐私与安全性}。
\item 页面下方的\ui{安全性}一栏会显示\ui{已阻止"BiWrite"以保护 Mac}，点旁边的\ui{仍要打开}。
\item 在弹出的确认框里再点\ui{仍要打开}，按要求输入登录密码或使用触控 ID。
\end{enumerate}
完成一次之后，以后可以直接打开。熟悉终端的用户也可以执行下面这条命令，效果相同：
\begin{center}\texttt{xattr -dr com.apple.quarantine /Applications/BiWrite.app}\end{center}

\subsection{Windows}
下载 \texttt{BiWrite\_\textless 版本\textgreater\_x64-setup.exe}，适用于 64 位 Windows 10 和 Windows 11。安装包同样没有受信任的代码签名，运行时 SmartScreen 会弹出\ui{Windows 已保护你的电脑}。点\ui{更多信息}，再点\ui{仍要运行}，之后按安装向导完成安装。

\subsection{安装 TeX 发行版}
编译 LaTeX 论文需要本机装有 TeX 发行版。macOS 安装 MacTeX，Windows 安装 TeX Live 或 MiKTeX。BiWrite 会在常见的安装位置自动查找，找到的发行版显示在\ui{设置}的\ui{LaTeX}一节（\fig{i-latex} ①）。装在别处时点\ui{指定文件夹…}（②），选择包含 \texttt{pdflatex} 的文件夹。③控制保存后是否自动编译。

\shot{i-latex}{①找到的 TeX 发行版，②指定文件夹，③保存后自动编译}

\section{配置模型接口}
\label{sec:provider}

\subsection{支持的接口}
翻译和写作助手都通过大模型接口完成。BiWrite 支持以下几类接口。
\begin{itemize}
\item OpenAI 兼容的 Chat Completions 接口，覆盖 OpenAI、DeepSeek、通义千问、Kimi、OpenRouter，以及 Ollama、LM Studio 等本地服务。
\item OpenAI Responses 接口，适用于需要推理强度和服务等级参数的推理模型。
\item Anthropic Messages 接口。
\end{itemize}
内置的\ui{模拟接口}不联网，只把每段文字倒序输出，用于在没有密钥时试用界面。

\subsection{添加接口}
\begin{enumerate}
\item 打开\ui{设置}（\cmdkey ,）。①是已添加的接口列表，单选按钮选中的接口负责翻译（\fig{c-add}）。
\item 在②的下拉框选择\ui{+ 添加接口…}，再选一个预设。
\item 在接口表单（\fig{c-form}）中核对 Base URL、模型名和接口类型（①），按需要选择推理强度和服务等级（②）。
\item 在密钥框粘贴密钥，点\ui{添加接口}。密钥保存在系统钥匙串（Windows 为凭据管理器）中，之后界面不再显示密钥原文。已有密钥时用\ui{管理号池…}（③）管理。
\item 点\ui{测试}（④）翻译一句话，确认接口可用。
\end{enumerate}

\shot{c-add}{①接口列表，②添加接口}
\shot{c-form}{接口表单：①接口类型与温度，②推理强度与服务等级，③管理号池，④测试}

\subsection{以第三方接口 AnyRouter 为例}
AnyRouter 是一个兼容 OpenAI Responses 接口的第三方服务。在预设中选择\ui{AnyRouter (GPT-6 Astra)}，预设会填好 Responses 接口、地址 \texttt{https://anyrouter.top/v1}、模型 \texttt{gpt-6-astra}、推理强度\ui{高}、服务等级\ui{Priority}、每个密钥并发 2、失败重试 10 次。粘贴一个或多个密钥后保存即可。其他第三方接口的配置方法相同，按服务商文档填写地址、模型和接口类型。请求日志会显示服务端实际返回的推理强度和服务等级，可以用来核对第三方服务是否按要求执行。

\subsection{号池：多个密钥一起用}
一个接口可以保存多个密钥，组成号池。在接口表单中点\ui{管理号池…}打开号池窗口（\fig{c-keys}）。
\begin{itemize}
\item ①是号池中的密钥。每个密钥可以起一个名字，表格只显示密钥末尾几位、状态和进行中的请求数。每个请求交给当前最空闲的可用密钥，某个密钥被限流时进入冷却，请求自动换到其他密钥，被拒绝的密钥暂停使用。
\item ②设置每个密钥的并发和失败重试次数。AnyRouter 每个账号同时只允许两个对话，所以设为 2。
\item ③是粘贴框：每行一个密钥，或者一行名称、下一行密钥，也可以写成\ui{名称: 密钥}，点\ui{追加密钥}加入号池。
\item ④\ui{从文件导入…}和\ui{导出到文件…}：导出把名称和密钥保存成文本文件，方便换电脑或分享给同事，导入读取这种文件并追加到号池。密钥从钥匙串直接写入文件，不经过界面。文件里是明文密钥，只发给你信任的人。
\end{itemize}

\shot[0.8]{c-keys}{号池：①密钥列表，②并发与重试，③粘贴新密钥，④从文件导入、导出到文件}

\subsection{翻译和写作分别用不同接口}
\ui{设置}中的\ui{写作助手}一节（\fig{c-assistant} ①）可以单独选择润色、修改和提问所用的模型。常见做法是用便宜快速的模型翻译，用更强的模型写作。选\ui{与翻译相同}时，两者使用同一个接口。

\shot{c-assistant}{①写作助手所用的模型}

\subsection{请求节奏}
\ui{请求节奏}一节（\fig{c-pacing}）控制发往接口的请求量。
\begin{itemize}
\item ①每次请求的段落数：默认 1，即一段一个请求。调到 2 到 8 时，多个新段落合并成一次请求，请求数随之减少，适合有频率限制的接口。改动过的段落始终单独发送，以便按原译文做最小修改。
\item ②并行请求数：同时进行的请求上限。
\item ③勾选\ui{按号池自动设置}时，并行数等于密钥个数乘以每个密钥的并发。
\end{itemize}

\shot{c-pacing}{①每次请求的段落数，②并行请求数，③按号池自动设置}

\subsection{界面语言}
\ui{设置}最上方的\ui{界面语言}（\fig{c-language} ①）可以在 English 和中文之间切换，只改变界面文字，不影响翻译方向。

\shot{c-language}{①界面语言}

\subsection{关闭窗口与托盘}
\label{sec:tray}
\ui{设置}中的\ui{窗口}一节（\fig{s-tray} ①）默认勾选\ui{关闭窗口时隐藏到托盘}。关闭窗口后，BiWrite 带着打开的文档继续运行，翻译和写作助手的任务照常进行。macOS 上点菜单栏中的 BiWrite 图标，Windows 上点任务栏通知区域的图标，可以选择\ui{显示 BiWrite}或\ui{退出 BiWrite}。macOS 上点程序坞中的图标也会重新显示窗口。退出时如果有未保存的改动，软件先询问。取消勾选后，关闭窗口即退出软件。

\shot{s-tray}{①关闭窗口时隐藏到托盘}

\section{快捷键}
\begin{center}
\begin{tabularx}{0.82\linewidth}{@{}lX@{}}
\toprule
macOS（Windows 用 Ctrl 代替 \cmdkey，Alt 代替 \optkey） & 作用 \\
\midrule
\cmdkey O & 打开文件 \\
\cmdkey S / \cmdkey\shiftkey S & 保存 / 另存为 \\
\cmdkey B & 编译当前 PDF \\
\cmdkey\optkey J & 在 PDF 中定位光标 \\
\cmdkey J & 打开或关闭写作助手 \\
\cmdkey\shiftkey P & 润色光标所在段落 \\
\cmdkey K & 聚焦写作助手的输入框 \\
\cmdkey , & 打开设置 \\
\cmdkey\shiftkey L & 打开请求日志 \\
\cmdkey F & 在源码中查找和替换 \\
\cmdkey Z & 撤销（包括采纳的修订） \\
\bottomrule
\end{tabularx}
\end{center}

\section{常见问题}

\paragraph{提示没有找到 TeX 发行版。}安装 MacTeX、TeX Live 或 MiKTeX 后，在 PDF 页签点\ui{重新检测}。装在非默认位置时，在\ui{设置}的\ui{LaTeX}一节指定文件夹。

\paragraph{编译失败或 PDF 不是最新的。}看 PDF 上方的问题列表，点击错误跳到源码行。没有日志信息的失败可以展开\ui{控制台输出}查看工具输出。编译超过 5 分钟会自动停止。

\paragraph{点击 PDF 没有反应。}由 BiWrite 编译出的 PDF 都带有 SyncTeX 信息。点在图片或空白处时，菜单会提示没有对应的源码。

\paragraph{点击 PDF 时提示正在重新编译。}导入或取消配对之后，以及未配对时切换语言之后，之前编译的 PDF 不再对应左侧的文字，PDF 上方会注明它编译于配对或切换语言之前。此时点击或定位光标会先重新编译这份 PDF，不会自动保存文件，编译完成后再点一次即可。

\paragraph{切换语言按钮一直在等待。}说明还有段落在翻译或翻译失败。翻译失败的段落可以在右侧点\ui{重试}，或用\ui{全部}菜单的\ui{继续翻译剩余段落}重试，也可以再点一次切换按钮，立即切换。

\paragraph{打开中文文件后右侧还是中文。}说明软件把原文判断成了英文，常见于中英文混合的文档。在\ui{全部}菜单中选\ui{把原文当作中文，翻译成英文}即可，文件内容不变。

\paragraph{插入的图表编译报错。}看结果卡片中列出的宏包，确认主文件导言区已加载。当前文件是主文件时，插入会自动加入这些宏包。

\paragraph{关闭窗口后软件去哪了。}关闭窗口默认隐藏到托盘，软件仍在运行。点菜单栏或任务栏通知区域的 BiWrite 图标即可回到窗口，或在\ui{设置}的\ui{窗口}一节关闭这一行为，见第~\ref{sec:tray}~节。

\paragraph{第三方接口经常限流。}在号池中加入多个密钥，把\ui{每个密钥的并发}设成服务商允许的数值，把\ui{每次请求的段落数}调大，并适当提高\ui{失败重试次数}。请求日志会显示每次限流和换密钥的经过。

\paragraph{导入中文文件时提示段落对不上。}两个文件中能对上的段落少于一半时不会配对。检查选的是否是同一篇文档的两种语言版本。

\paragraph{数据保存在哪里。}翻译缓存位于应用数据文件夹中的 \texttt{cache.sqlite3}，可以在\ui{设置}的\ui{翻译缓存}一节查看大小和清空。密钥只保存在系统钥匙串中。

\end{document}
